\documentclass[aspectratio=169,10pt]{beamer}
\makeatletter\def\input@path{{../../estilo/}{../estilo/}{estilo/}}\makeatother
\usepackage{devmobile}
\modulo[Kotlin]{01}{Ambiente de Desenvolvimento}{Android Studio, SDK, emulador, celular real, estrutura de projeto, Gradle e o primeiro app}

\begin{document}
\capa

\begin{frame}{Objetivos do módulo}
  \begin{conceito}[Ao final deste módulo você será capaz de]
    \begin{itemize}
      \item Instalar e configurar o Android Studio e o Android SDK.
      \item Executar apps no emulador e em um celular real (USB e Wi-Fi).
      \item Navegar pela estrutura de um projeto Android e entender cada arquivo.
      \item Entender o papel do Gradle, do \emph{version catalog} e das variantes de build.
      \item Criar, executar e depurar o primeiro app com Jetpack Compose.
      \item Versionar o projeto com Git e GitHub.
    \end{itemize}
  \end{conceito}
\end{frame}

\sumario

% =============================================================================
%  Fundamentos compartilhados --- Módulo 01: ferramentas de desenvolvimento
%  Incluído pelas trilhas Kotlin e Flutter com % =============================================================================
%  Fundamentos compartilhados --- Módulo 01: ferramentas de desenvolvimento
%  Incluído pelas trilhas Kotlin e Flutter com % =============================================================================
%  Fundamentos compartilhados --- Módulo 01: ferramentas de desenvolvimento
%  Incluído pelas trilhas Kotlin e Flutter com \input{../fundamentos/f01-ambiente}
% =============================================================================
\section{Fundamentos: ferramentas de desenvolvimento}

\begin{frame}{Do código-fonte ao programa}
  \centering
  \begin{tikzpicture}[node distance=5mm,every node/.style={font=\scriptsize,align=center,minimum height=10mm}]
    \node[caixa,text width=1.9cm] (src) {Código-fonte\\ \tiny texto legível};
    \node[caixapetroleo,right=of src,text width=1.9cm] (comp) {Compilador\\ \tiny verifica e traduz};
    \node[caixaambar,right=of comp,text width=2.2cm] (bc) {\emph{Bytecode} ou\\ código de máquina};
    \node[caixapetroleo,right=of bc,text width=1.9cm] (pk) {Empacotador\\ \tiny junta recursos};
    \node[caixaescura,right=of pk,text width=1.9cm] (app) {App\\ instalável};
    \foreach \a/\b in {src/comp,comp/bc,bc/pk,pk/app} \draw[seta] (\a) -- (\b);
  \end{tikzpicture}
  \vspace{4pt}
  \begin{columns}[T]
    \begin{column}{0.48\textwidth}
      \begin{definicao}[Compilador e interpretador]
        O \textbf{compilador} traduz o programa inteiro antes da execução e aponta erros cedo. O \textbf{interpretador} lê e executa o código aos poucos, durante a execução.
      \end{definicao}
      \begin{definicao}[Máquina virtual]
        Programa que executa um \emph{bytecode} intermediário em qualquer processador: a JVM e o ART (Kotlin) e a VM do Dart.
      \end{definicao}
    \end{column}
    \begin{column}{0.48\textwidth}
      \begin{definicao}[JIT e AOT]
        \textbf{JIT}: compila enquanto o programa roda; permite trocar código em execução (\emph{hot reload}).
        \textbf{AOT}: compila tudo antes de distribuir; inicia rápido e roda com desempenho previsível.
      \end{definicao}
      \begin{dica}
        As duas trilhas usam os dois: JIT para desenvolver com agilidade, AOT para o app que chega ao usuário (Módulo 15).
      \end{dica}
    \end{column}
  \end{columns}
\end{frame}

\begin{frame}{SDK, IDE e cadeia de ferramentas}
  \begin{columns}[T]
    \begin{column}{0.48\textwidth}
      \begin{definicao}[SDK --- \emph{Software Development Kit}]
        Pacote com bibliotecas, compiladores, ferramentas de linha de comando, documentação e imagens de emulador para desenvolver para uma plataforma.
      \end{definicao}
      \begin{definicao}[IDE --- ambiente integrado]
        Editor que reúne escrita de código, compilação, execução, depuração, análise e integração com Git.
      \end{definicao}
      \begin{analogia}
        O SDK é a caixa de ferramentas; a IDE é a bancada organizada onde você trabalha com elas.
      \end{analogia}
    \end{column}
    \begin{column}{0.5\textwidth}
      \scriptsize
      \begin{tabularx}{\linewidth}{@{}l L L@{}}
        \toprule
        & \cabecalho{Trilha Kotlin} & \cabecalho{Trilha Flutter} \\
        \midrule
        Linguagem & Kotlin (sobre o JDK) & Dart (Dart SDK) \\
        SDK da plataforma & Android SDK & Flutter SDK + Android SDK + Xcode \\
        IDE & Android Studio & VS Code ou Android Studio \\
        \emph{Build} & Gradle & \cod{flutter} CLI (usa Gradle e Xcode) \\
        Aparelho & ADB & ADB e \cod{flutter devices} \\
        Pacotes & Maven Central, Google Maven & pub.dev \\
        \bottomrule
      \end{tabularx}
    \end{column}
  \end{columns}
\end{frame}

\begin{frame}{Sistemas de \emph{build} e dependências}
  \begin{columns}[c]
    \begin{column}{0.46\textwidth}
      \centering
      \begin{tikzpicture}[node distance=4mm,every node/.style={font=\tiny,minimum width=1.6cm}]
        \node[caixaescura] (apk) {gerar pacote};
        \node[caixapetroleo,below left=5mm and -2mm of apk] (comp) {compilar};
        \node[caixapetroleo,below right=5mm and -2mm of apk] (res) {processar recursos};
        \node[caixa,below=5mm of comp] (dep) {baixar dependências};
        \node[caixa,below=5mm of res] (gen) {gerar código};
        \draw[seta] (comp) -- (apk); \draw[seta] (res) -- (apk);
        \draw[seta] (dep) -- (comp); \draw[seta] (gen) -- (comp);
      \end{tikzpicture}\\[2pt]
      {\tiny Grafo de tarefas: cada uma só roda depois das que ela usa.}
    \end{column}
    \begin{column}{0.52\textwidth}
      \begin{definicao}[Sistema de \emph{build}]
        Ferramenta que transforma o projeto em app executando \textbf{tarefas} organizadas num grafo de dependências. Só refaz o que mudou (\emph{build} incremental) e guarda resultados em cache.
      \end{definicao}
      \begin{definicao}[Dependência transitiva]
        A biblioteca que você usa depende de outras, que dependem de outras. O gerenciador resolve a árvore inteira e escolhe versões compatíveis.
      \end{definicao}
      \begin{atencao}
        Fixe as versões com arquivos de trava (catálogo de versões, \cod{pubspec.lock}) para que todos na equipe compilem o mesmo app.
      \end{atencao}
    \end{column}
  \end{columns}
\end{frame}

\begin{frame}{Emuladores, simuladores e aparelhos reais}
  \begin{columns}[T]
    \begin{column}{0.5\textwidth}
      \scriptsize
      \begin{tabularx}{\linewidth}{@{}l L L@{}}
        \toprule
        & \cabecalho{Vantagem} & \cabecalho{Limite} \\
        \midrule
        Emulador Android & vários tamanhos e versões, GPS e sensores simulados & sem Bluetooth, NFC e USB reais \\
        Simulador iOS & rápido, integrado ao Xcode & roda código do Mac; não imita o hardware \\
        Aparelho real & desempenho e hardware verdadeiros & cada modelo exige um aparelho \\
        \bottomrule
      \end{tabularx}
    \end{column}
    \begin{column}{0.48\textwidth}
      \begin{definicao}[Emulador e simulador]
        O \textbf{emulador} imita o aparelho, inclusive o processador e o sistema completo. O \textbf{simulador} apenas reproduz o comportamento do sistema, executando o app compilado para o computador.
      \end{definicao}
      \begin{pordentro}[Por que o emulador exige virtualização]
        O emulador Android roda uma imagem completa do sistema numa máquina virtual. Com aceleração por hardware (Intel VT-x, AMD-V, KVM no Linux, Hypervisor no macOS), as instruções rodam direto no processador e o emulador fica rápido.
      \end{pordentro}
    \end{column}
  \end{columns}
\end{frame}

\begin{frame}{Depuração}
  \begin{columns}[T]
    \begin{column}{0.5\textwidth}
      \begin{definicao}[Depurador]
        Ferramenta que pausa o programa num ponto escolhido (\emph{breakpoint}) para inspecionar variáveis, a pilha de chamadas e executar passo a passo.
      \end{definicao}
      \scriptsize
      \begin{tabularx}{\linewidth}{@{}l L@{}}
        \toprule
        \cabecalho{Comando} & \cabecalho{O que faz} \\
        \midrule
        \emph{Step over} & executa a linha e para na próxima \\
        \emph{Step into} & entra na função chamada \\
        \emph{Step out} & termina a função atual e volta \\
        \emph{Resume} & continua até o próximo \emph{breakpoint} \\
        \bottomrule
      \end{tabularx}
    \end{column}
    \begin{column}{0.48\textwidth}
      \begin{conceito}[Método científico da depuração]
        \begin{enumerate}
          \item Reproduza o problema de forma confiável.
          \item Formule uma hipótese sobre a causa.
          \item Teste a hipótese (\emph{breakpoint}, \emph{log}, teste automatizado).
          \item Corrija e escreva um teste que impeça a volta do defeito.
        \end{enumerate}
      \end{conceito}
      \begin{dica}
        Leia a mensagem de erro e a pilha de chamadas inteiras: a primeira linha do \textbf{seu} código costuma apontar a causa.
      \end{dica}
    \end{column}
  \end{columns}
\end{frame}

% =============================================================================
\section{Fundamentos: ferramentas de desenvolvimento}

\begin{frame}{Do código-fonte ao programa}
  \centering
  \begin{tikzpicture}[node distance=5mm,every node/.style={font=\scriptsize,align=center,minimum height=10mm}]
    \node[caixa,text width=1.9cm] (src) {Código-fonte\\ \tiny texto legível};
    \node[caixapetroleo,right=of src,text width=1.9cm] (comp) {Compilador\\ \tiny verifica e traduz};
    \node[caixaambar,right=of comp,text width=2.2cm] (bc) {\emph{Bytecode} ou\\ código de máquina};
    \node[caixapetroleo,right=of bc,text width=1.9cm] (pk) {Empacotador\\ \tiny junta recursos};
    \node[caixaescura,right=of pk,text width=1.9cm] (app) {App\\ instalável};
    \foreach \a/\b in {src/comp,comp/bc,bc/pk,pk/app} \draw[seta] (\a) -- (\b);
  \end{tikzpicture}
  \vspace{4pt}
  \begin{columns}[T]
    \begin{column}{0.48\textwidth}
      \begin{definicao}[Compilador e interpretador]
        O \textbf{compilador} traduz o programa inteiro antes da execução e aponta erros cedo. O \textbf{interpretador} lê e executa o código aos poucos, durante a execução.
      \end{definicao}
      \begin{definicao}[Máquina virtual]
        Programa que executa um \emph{bytecode} intermediário em qualquer processador: a JVM e o ART (Kotlin) e a VM do Dart.
      \end{definicao}
    \end{column}
    \begin{column}{0.48\textwidth}
      \begin{definicao}[JIT e AOT]
        \textbf{JIT}: compila enquanto o programa roda; permite trocar código em execução (\emph{hot reload}).
        \textbf{AOT}: compila tudo antes de distribuir; inicia rápido e roda com desempenho previsível.
      \end{definicao}
      \begin{dica}
        As duas trilhas usam os dois: JIT para desenvolver com agilidade, AOT para o app que chega ao usuário (Módulo 15).
      \end{dica}
    \end{column}
  \end{columns}
\end{frame}

\begin{frame}{SDK, IDE e cadeia de ferramentas}
  \begin{columns}[T]
    \begin{column}{0.48\textwidth}
      \begin{definicao}[SDK --- \emph{Software Development Kit}]
        Pacote com bibliotecas, compiladores, ferramentas de linha de comando, documentação e imagens de emulador para desenvolver para uma plataforma.
      \end{definicao}
      \begin{definicao}[IDE --- ambiente integrado]
        Editor que reúne escrita de código, compilação, execução, depuração, análise e integração com Git.
      \end{definicao}
      \begin{analogia}
        O SDK é a caixa de ferramentas; a IDE é a bancada organizada onde você trabalha com elas.
      \end{analogia}
    \end{column}
    \begin{column}{0.5\textwidth}
      \scriptsize
      \begin{tabularx}{\linewidth}{@{}l L L@{}}
        \toprule
        & \cabecalho{Trilha Kotlin} & \cabecalho{Trilha Flutter} \\
        \midrule
        Linguagem & Kotlin (sobre o JDK) & Dart (Dart SDK) \\
        SDK da plataforma & Android SDK & Flutter SDK + Android SDK + Xcode \\
        IDE & Android Studio & VS Code ou Android Studio \\
        \emph{Build} & Gradle & \cod{flutter} CLI (usa Gradle e Xcode) \\
        Aparelho & ADB & ADB e \cod{flutter devices} \\
        Pacotes & Maven Central, Google Maven & pub.dev \\
        \bottomrule
      \end{tabularx}
    \end{column}
  \end{columns}
\end{frame}

\begin{frame}{Sistemas de \emph{build} e dependências}
  \begin{columns}[c]
    \begin{column}{0.46\textwidth}
      \centering
      \begin{tikzpicture}[node distance=4mm,every node/.style={font=\tiny,minimum width=1.6cm}]
        \node[caixaescura] (apk) {gerar pacote};
        \node[caixapetroleo,below left=5mm and -2mm of apk] (comp) {compilar};
        \node[caixapetroleo,below right=5mm and -2mm of apk] (res) {processar recursos};
        \node[caixa,below=5mm of comp] (dep) {baixar dependências};
        \node[caixa,below=5mm of res] (gen) {gerar código};
        \draw[seta] (comp) -- (apk); \draw[seta] (res) -- (apk);
        \draw[seta] (dep) -- (comp); \draw[seta] (gen) -- (comp);
      \end{tikzpicture}\\[2pt]
      {\tiny Grafo de tarefas: cada uma só roda depois das que ela usa.}
    \end{column}
    \begin{column}{0.52\textwidth}
      \begin{definicao}[Sistema de \emph{build}]
        Ferramenta que transforma o projeto em app executando \textbf{tarefas} organizadas num grafo de dependências. Só refaz o que mudou (\emph{build} incremental) e guarda resultados em cache.
      \end{definicao}
      \begin{definicao}[Dependência transitiva]
        A biblioteca que você usa depende de outras, que dependem de outras. O gerenciador resolve a árvore inteira e escolhe versões compatíveis.
      \end{definicao}
      \begin{atencao}
        Fixe as versões com arquivos de trava (catálogo de versões, \cod{pubspec.lock}) para que todos na equipe compilem o mesmo app.
      \end{atencao}
    \end{column}
  \end{columns}
\end{frame}

\begin{frame}{Emuladores, simuladores e aparelhos reais}
  \begin{columns}[T]
    \begin{column}{0.5\textwidth}
      \scriptsize
      \begin{tabularx}{\linewidth}{@{}l L L@{}}
        \toprule
        & \cabecalho{Vantagem} & \cabecalho{Limite} \\
        \midrule
        Emulador Android & vários tamanhos e versões, GPS e sensores simulados & sem Bluetooth, NFC e USB reais \\
        Simulador iOS & rápido, integrado ao Xcode & roda código do Mac; não imita o hardware \\
        Aparelho real & desempenho e hardware verdadeiros & cada modelo exige um aparelho \\
        \bottomrule
      \end{tabularx}
    \end{column}
    \begin{column}{0.48\textwidth}
      \begin{definicao}[Emulador e simulador]
        O \textbf{emulador} imita o aparelho, inclusive o processador e o sistema completo. O \textbf{simulador} apenas reproduz o comportamento do sistema, executando o app compilado para o computador.
      \end{definicao}
      \begin{pordentro}[Por que o emulador exige virtualização]
        O emulador Android roda uma imagem completa do sistema numa máquina virtual. Com aceleração por hardware (Intel VT-x, AMD-V, KVM no Linux, Hypervisor no macOS), as instruções rodam direto no processador e o emulador fica rápido.
      \end{pordentro}
    \end{column}
  \end{columns}
\end{frame}

\begin{frame}{Depuração}
  \begin{columns}[T]
    \begin{column}{0.5\textwidth}
      \begin{definicao}[Depurador]
        Ferramenta que pausa o programa num ponto escolhido (\emph{breakpoint}) para inspecionar variáveis, a pilha de chamadas e executar passo a passo.
      \end{definicao}
      \scriptsize
      \begin{tabularx}{\linewidth}{@{}l L@{}}
        \toprule
        \cabecalho{Comando} & \cabecalho{O que faz} \\
        \midrule
        \emph{Step over} & executa a linha e para na próxima \\
        \emph{Step into} & entra na função chamada \\
        \emph{Step out} & termina a função atual e volta \\
        \emph{Resume} & continua até o próximo \emph{breakpoint} \\
        \bottomrule
      \end{tabularx}
    \end{column}
    \begin{column}{0.48\textwidth}
      \begin{conceito}[Método científico da depuração]
        \begin{enumerate}
          \item Reproduza o problema de forma confiável.
          \item Formule uma hipótese sobre a causa.
          \item Teste a hipótese (\emph{breakpoint}, \emph{log}, teste automatizado).
          \item Corrija e escreva um teste que impeça a volta do defeito.
        \end{enumerate}
      \end{conceito}
      \begin{dica}
        Leia a mensagem de erro e a pilha de chamadas inteiras: a primeira linha do \textbf{seu} código costuma apontar a causa.
      \end{dica}
    \end{column}
  \end{columns}
\end{frame}

% =============================================================================
\section{Fundamentos: ferramentas de desenvolvimento}

\begin{frame}{Do código-fonte ao programa}
  \centering
  \begin{tikzpicture}[node distance=5mm,every node/.style={font=\scriptsize,align=center,minimum height=10mm}]
    \node[caixa,text width=1.9cm] (src) {Código-fonte\\ \tiny texto legível};
    \node[caixapetroleo,right=of src,text width=1.9cm] (comp) {Compilador\\ \tiny verifica e traduz};
    \node[caixaambar,right=of comp,text width=2.2cm] (bc) {\emph{Bytecode} ou\\ código de máquina};
    \node[caixapetroleo,right=of bc,text width=1.9cm] (pk) {Empacotador\\ \tiny junta recursos};
    \node[caixaescura,right=of pk,text width=1.9cm] (app) {App\\ instalável};
    \foreach \a/\b in {src/comp,comp/bc,bc/pk,pk/app} \draw[seta] (\a) -- (\b);
  \end{tikzpicture}
  \vspace{4pt}
  \begin{columns}[T]
    \begin{column}{0.48\textwidth}
      \begin{definicao}[Compilador e interpretador]
        O \textbf{compilador} traduz o programa inteiro antes da execução e aponta erros cedo. O \textbf{interpretador} lê e executa o código aos poucos, durante a execução.
      \end{definicao}
      \begin{definicao}[Máquina virtual]
        Programa que executa um \emph{bytecode} intermediário em qualquer processador: a JVM e o ART (Kotlin) e a VM do Dart.
      \end{definicao}
    \end{column}
    \begin{column}{0.48\textwidth}
      \begin{definicao}[JIT e AOT]
        \textbf{JIT}: compila enquanto o programa roda; permite trocar código em execução (\emph{hot reload}).
        \textbf{AOT}: compila tudo antes de distribuir; inicia rápido e roda com desempenho previsível.
      \end{definicao}
      \begin{dica}
        As duas trilhas usam os dois: JIT para desenvolver com agilidade, AOT para o app que chega ao usuário (Módulo 15).
      \end{dica}
    \end{column}
  \end{columns}
\end{frame}

\begin{frame}{SDK, IDE e cadeia de ferramentas}
  \begin{columns}[T]
    \begin{column}{0.48\textwidth}
      \begin{definicao}[SDK --- \emph{Software Development Kit}]
        Pacote com bibliotecas, compiladores, ferramentas de linha de comando, documentação e imagens de emulador para desenvolver para uma plataforma.
      \end{definicao}
      \begin{definicao}[IDE --- ambiente integrado]
        Editor que reúne escrita de código, compilação, execução, depuração, análise e integração com Git.
      \end{definicao}
      \begin{analogia}
        O SDK é a caixa de ferramentas; a IDE é a bancada organizada onde você trabalha com elas.
      \end{analogia}
    \end{column}
    \begin{column}{0.5\textwidth}
      \scriptsize
      \begin{tabularx}{\linewidth}{@{}l L L@{}}
        \toprule
        & \cabecalho{Trilha Kotlin} & \cabecalho{Trilha Flutter} \\
        \midrule
        Linguagem & Kotlin (sobre o JDK) & Dart (Dart SDK) \\
        SDK da plataforma & Android SDK & Flutter SDK + Android SDK + Xcode \\
        IDE & Android Studio & VS Code ou Android Studio \\
        \emph{Build} & Gradle & \cod{flutter} CLI (usa Gradle e Xcode) \\
        Aparelho & ADB & ADB e \cod{flutter devices} \\
        Pacotes & Maven Central, Google Maven & pub.dev \\
        \bottomrule
      \end{tabularx}
    \end{column}
  \end{columns}
\end{frame}

\begin{frame}{Sistemas de \emph{build} e dependências}
  \begin{columns}[c]
    \begin{column}{0.46\textwidth}
      \centering
      \begin{tikzpicture}[node distance=4mm,every node/.style={font=\tiny,minimum width=1.6cm}]
        \node[caixaescura] (apk) {gerar pacote};
        \node[caixapetroleo,below left=5mm and -2mm of apk] (comp) {compilar};
        \node[caixapetroleo,below right=5mm and -2mm of apk] (res) {processar recursos};
        \node[caixa,below=5mm of comp] (dep) {baixar dependências};
        \node[caixa,below=5mm of res] (gen) {gerar código};
        \draw[seta] (comp) -- (apk); \draw[seta] (res) -- (apk);
        \draw[seta] (dep) -- (comp); \draw[seta] (gen) -- (comp);
      \end{tikzpicture}\\[2pt]
      {\tiny Grafo de tarefas: cada uma só roda depois das que ela usa.}
    \end{column}
    \begin{column}{0.52\textwidth}
      \begin{definicao}[Sistema de \emph{build}]
        Ferramenta que transforma o projeto em app executando \textbf{tarefas} organizadas num grafo de dependências. Só refaz o que mudou (\emph{build} incremental) e guarda resultados em cache.
      \end{definicao}
      \begin{definicao}[Dependência transitiva]
        A biblioteca que você usa depende de outras, que dependem de outras. O gerenciador resolve a árvore inteira e escolhe versões compatíveis.
      \end{definicao}
      \begin{atencao}
        Fixe as versões com arquivos de trava (catálogo de versões, \cod{pubspec.lock}) para que todos na equipe compilem o mesmo app.
      \end{atencao}
    \end{column}
  \end{columns}
\end{frame}

\begin{frame}{Emuladores, simuladores e aparelhos reais}
  \begin{columns}[T]
    \begin{column}{0.5\textwidth}
      \scriptsize
      \begin{tabularx}{\linewidth}{@{}l L L@{}}
        \toprule
        & \cabecalho{Vantagem} & \cabecalho{Limite} \\
        \midrule
        Emulador Android & vários tamanhos e versões, GPS e sensores simulados & sem Bluetooth, NFC e USB reais \\
        Simulador iOS & rápido, integrado ao Xcode & roda código do Mac; não imita o hardware \\
        Aparelho real & desempenho e hardware verdadeiros & cada modelo exige um aparelho \\
        \bottomrule
      \end{tabularx}
    \end{column}
    \begin{column}{0.48\textwidth}
      \begin{definicao}[Emulador e simulador]
        O \textbf{emulador} imita o aparelho, inclusive o processador e o sistema completo. O \textbf{simulador} apenas reproduz o comportamento do sistema, executando o app compilado para o computador.
      \end{definicao}
      \begin{pordentro}[Por que o emulador exige virtualização]
        O emulador Android roda uma imagem completa do sistema numa máquina virtual. Com aceleração por hardware (Intel VT-x, AMD-V, KVM no Linux, Hypervisor no macOS), as instruções rodam direto no processador e o emulador fica rápido.
      \end{pordentro}
    \end{column}
  \end{columns}
\end{frame}

\begin{frame}{Depuração}
  \begin{columns}[T]
    \begin{column}{0.5\textwidth}
      \begin{definicao}[Depurador]
        Ferramenta que pausa o programa num ponto escolhido (\emph{breakpoint}) para inspecionar variáveis, a pilha de chamadas e executar passo a passo.
      \end{definicao}
      \scriptsize
      \begin{tabularx}{\linewidth}{@{}l L@{}}
        \toprule
        \cabecalho{Comando} & \cabecalho{O que faz} \\
        \midrule
        \emph{Step over} & executa a linha e para na próxima \\
        \emph{Step into} & entra na função chamada \\
        \emph{Step out} & termina a função atual e volta \\
        \emph{Resume} & continua até o próximo \emph{breakpoint} \\
        \bottomrule
      \end{tabularx}
    \end{column}
    \begin{column}{0.48\textwidth}
      \begin{conceito}[Método científico da depuração]
        \begin{enumerate}
          \item Reproduza o problema de forma confiável.
          \item Formule uma hipótese sobre a causa.
          \item Teste a hipótese (\emph{breakpoint}, \emph{log}, teste automatizado).
          \item Corrija e escreva um teste que impeça a volta do defeito.
        \end{enumerate}
      \end{conceito}
      \begin{dica}
        Leia a mensagem de erro e a pilha de chamadas inteiras: a primeira linha do \textbf{seu} código costuma apontar a causa.
      \end{dica}
    \end{column}
  \end{columns}
\end{frame}


% =============================================================================
\section{Ferramentas}

\begin{frame}{Android Studio}
  \begin{columns}[T]
    \begin{column}{0.55\textwidth}
      \begin{itemize}
        \item IDE oficial do Android, baseada no IntelliJ IDEA (JetBrains).
        \item Gratuita, para Windows, macOS, Linux e ChromeOS.
        \item Inclui o JDK, o SDK Manager, o emulador, o \emph{profiler} e o inspetor de layout.
        \item Assistente de IA integrado (Gemini).
      \end{itemize}
      \begin{atencao}[Requisitos práticos]
        \begin{itemize}
          \item \textbf{RAM:} 16\,GB; com 8\,GB, prefira o celular real.
          \item \textbf{Disco:} SSD com 20\,GB livres.
          \item \textbf{CPU:} virtualização ativada na BIOS.
        \end{itemize}
      \end{atencao}
    \end{column}
    \begin{column}{0.42\textwidth}
      \begin{tikzpicture}[node distance=1.5mm]
        \node[caixaescura,minimum width=4.6cm] (as) {\textbf{Android Studio}};
        \node[caixa,minimum width=2.25cm,below=of as.south west,anchor=north west] (ed) {Editor\\ + Compose Preview};
        \node[caixa,minimum width=2.25cm,below=of as.south east,anchor=north east] (gr) {Gradle\\ (build)};
        \node[caixa,minimum width=2.25cm,below=of ed] (sdk) {SDK Manager};
        \node[caixa,minimum width=2.25cm,below=of gr] (avd) {Device Manager\\ (emulador)};
        \node[caixaambar,minimum width=2.25cm,below=of sdk] (lc) {Logcat};
        \node[caixaambar,minimum width=2.25cm,below=of avd] (pf) {Profiler e\\ App Inspection};
      \end{tikzpicture}
    \end{column}
  \end{columns}
\end{frame}

\begin{frame}{Instalação passo a passo}
  \begin{enumerate}
    \item Baixe o instalador em \url{https://developer.android.com/studio}.
    \item \textbf{Windows:} execute o \arq{.exe}. \textbf{Linux:} extraia o \arq{.tar.gz} em \arq{/opt} e rode \arq{bin/studio.sh}. \textbf{macOS:} arraste para Aplicativos.
    \item No assistente de primeira execução escolha \textbf{Standard}: ele baixa o SDK, as ferramentas de build e uma imagem do emulador.
    \item Abra \textbf{Settings $\rightarrow$ Languages \& Frameworks $\rightarrow$ Android SDK} e confira os pacotes.
  \end{enumerate}
  \vspace{2pt}
  \begin{columns}[T]
    \begin{column}{0.48\textwidth}
      \begin{conceito}[Pacotes essenciais do SDK]
        \begin{itemize}
          \item \textbf{SDK Platform} da API alvo (ex.: Android 16).
          \item \textbf{Build-Tools} e \textbf{Platform-Tools} (contém o \cod{adb}).
          \item \textbf{Android Emulator} e uma \emph{system image}.
        \end{itemize}
      \end{conceito}
    \end{column}
    \begin{column}{0.48\textwidth}
      \begin{dica}[Linux: aceleração do emulador]
        Instale o KVM e adicione seu usuário ao grupo \cod{kvm}. Sem aceleração por hardware o emulador fica lento demais para uso.
      \end{dica}
    \end{column}
  \end{columns}
\end{frame}

% =============================================================================
\section{Executando o app}

\begin{frame}{Emulador (AVD)}
  \begin{columns}[T]
    \begin{column}{0.56\textwidth}
      \small
      \termo{AVD} (\emph{Android Virtual Device}) é uma configuração de aparelho virtual: modelo, tela, versão do Android e memória.
      \begin{enumerate}
        \item \textbf{Device Manager $\rightarrow$ +} (Create Virtual Device).
        \item Escolha um perfil (ex.: \emph{Pixel}, \emph{Medium Phone}).
        \item Escolha uma \emph{system image} com \textbf{Google Play} ou \textbf{Google APIs}.
        \item Prefira imagens \textbf{x86\_64} (PC) ou \textbf{arm64} (Mac Apple Silicon).
      \end{enumerate}
      \begin{dica}
        O emulador simula GPS, rotação, bateria, rede lenta e sensores em \emph{Extended Controls} (\textbf{\ldots}).
      \end{dica}
    \end{column}
    \begin{column}{0.40\textwidth}
      \centering
      \begin{tikzpicture}
        \node[caixaescura,minimum width=3.8cm] (pc) {Computador};
        \node[caixa,minimum width=3.8cm,below=2mm of pc] (hyp) {Hipervisor\\ (KVM, WHPX, HVF)};
        \node[caixaambar,minimum width=3.8cm,below=2mm of hyp] (emu) {Emulador QEMU};
        \node[below=2mm of emu] (cel) {\celular[0.45]{Android}};
        \draw[seta] (emu) -- (cel);
      \end{tikzpicture}
    \end{column}
  \end{columns}
\end{frame}

\begin{frame}{Celular real: depuração USB e Wi-Fi}
  \begin{columns}[T]
    \begin{column}{0.5\textwidth}
      \begin{conceito}[Ativando as opções do desenvolvedor]
        \begin{enumerate}
          \item \textbf{Configurações $\rightarrow$ Sobre o telefone}.
          \item Toque 7 vezes em \textbf{Número da versão}.
          \item Em \textbf{Opções do desenvolvedor}, ative \textbf{Depuração USB}.
          \item Conecte o cabo e aceite a impressão digital RSA no celular.
        \end{enumerate}
      \end{conceito}
    \end{column}
    \begin{column}{0.48\textwidth}
      \begin{conceito}[Depuração por Wi-Fi (Android 11+)]
        \begin{enumerate}
          \item Celular e computador na mesma rede.
          \item Ative \textbf{Depuração por Wi-Fi}.
          \item No Android Studio: \textbf{Device Manager $\rightarrow$ Pair devices using Wi-Fi} e leia o QR code.
        \end{enumerate}
      \end{conceito}
    \end{column}
  \end{columns}
  \begin{atencao}[Aparelho real é obrigatório para]
    Bluetooth/BLE, NFC, câmera de verdade, sensores físicos e medições realistas de desempenho e bateria. Projetos com ESP32 (Módulo 13) \textbf{exigem} celular real.
  \end{atencao}
\end{frame}

\begin{frame}[fragile]{ADB: \emph{Android Debug Bridge}}
  \begin{columns}[T]
    \begin{column}{0.56\textwidth}
\begin{terminal}[title=Comandos úteis]
adb devices                 # lista aparelhos conectados
adb install app-debug.apk   # instala um APK
adb uninstall com.exemplo.app
adb logcat -s LabMonitor    # filtra logs pela tag
adb shell                   # terminal no aparelho
adb shell pm list packages | grep exemplo
adb push foto.jpg /sdcard/Download/
adb pull /sdcard/Download/log.txt .
adb shell screencap -p /sdcard/tela.png
adb reverse tcp:8080 tcp:8080  # localhost do PC
\end{terminal}
    \end{column}
    \begin{column}{0.42\textwidth}
      \centering
      \begin{tikzpicture}[node distance=6mm]
        \node[caixaescura,minimum width=3.4cm] (cli) {cliente \cod{adb}\\ (terminal / IDE)};
        \node[caixa,minimum width=3.4cm,below=of cli] (srv) {servidor \cod{adb}\\ (porta 5037 no PC)};
        \node[caixaambar,minimum width=3.4cm,below=of srv] (d) {\cod{adbd}\\ (no aparelho)};
        \draw[seta] (cli) -- (srv);
        \draw[seta] (srv) -- node[rotulo,right]{USB ou Wi-Fi} (d);
      \end{tikzpicture}
      \begin{dica}
        \cod{adb reverse} permite que o app no celular acesse uma API rodando no seu computador.
      \end{dica}
    \end{column}
  \end{columns}
\end{frame}

% =============================================================================
\section{Estrutura de um projeto}

\begin{frame}{Criando o projeto}
  \begin{columns}[T]
    \begin{column}{0.5\textwidth}
      \begin{enumerate}
        \item \textbf{File $\rightarrow$ New $\rightarrow$ New Project}.
        \item Modelo \textbf{Empty Activity} (já vem com Compose).
        \item Preencha:
          \begin{itemize}
            \item \textbf{Name}: LabMonitor
            \item \textbf{Package name}: \cod{br.edu.ifma.labmonitor}
            \item \textbf{Minimum SDK}: API 26
            \item \textbf{Build configuration language}: Kotlin DSL
          \end{itemize}
        \item Aguarde o \textbf{Gradle Sync} terminar.
      \end{enumerate}
    \end{column}
    \begin{column}{0.48\textwidth}
      \begin{conceito}[Package name / applicationId]
        Identificador \textbf{único e permanente} do app na Google Play. Use o domínio invertido da instituição ou o seu. Depois de publicado, \textbf{não pode ser alterado}.
      \end{conceito}
      \begin{dica}
        Use a visualização \textbf{Android} no painel de projeto para ver a estrutura lógica, e \textbf{Project} para ver as pastas reais no disco.
      \end{dica}
    \end{column}
  \end{columns}
\end{frame}

\begin{frame}[fragile]{Árvore de diretórios}
  \begin{columns}[T]
    \begin{column}{0.46\textwidth}
\begin{codigo}[title=LabMonitor/,listing options={style=dm,language=bash,basicstyle=\ttfamily\tiny}]{bash}
LabMonitor/
|-- app/
|   |-- build.gradle.kts      (módulo)
|   |-- proguard-rules.pro
|   `-- src/
|       |-- main/
|       |   |-- AndroidManifest.xml
|       |   |-- java/.../labmonitor/
|       |   |   |-- MainActivity.kt
|       |   |   `-- ui/theme/
|       |   `-- res/
|       |       |-- drawable/  mipmap-*/
|       |       `-- values/strings.xml
|       |-- test/            (unitários)
|       `-- androidTest/     (instrumentados)
|-- gradle/
|   |-- libs.versions.toml   (catálogo)
|   `-- wrapper/
|-- build.gradle.kts          (projeto)
|-- settings.gradle.kts
|-- gradle.properties
`-- gradlew, gradlew.bat
\end{codigo}
    \end{column}
    \begin{column}{0.52\textwidth}
      \small
      \begin{itemize}
        \item \arq{app/}: o módulo do aplicativo. Projetos grandes têm vários módulos.
        \item \arq{src/main/java}: código Kotlin (o nome da pasta é histórico).
        \item \arq{res/}: recursos que não são código: imagens, textos, ícones, cores.
        \item \arq{test/}: testes que rodam na JVM do computador.
        \item \arq{androidTest/}: testes que rodam no aparelho.
        \item \arq{gradlew}: \emph{wrapper} que baixa a versão certa do Gradle. Sempre use-o, inclusive na integração contínua.
      \end{itemize}
      \begin{dica}
        Os textos da interface ficam em \arq{strings.xml}. Isso facilita a tradução do app.
      \end{dica}
    \end{column}
  \end{columns}
\end{frame}

\begin{frame}[fragile]{AndroidManifest.xml}
  \begin{columns}[T]
    \begin{column}{0.58\textwidth}
\begin{codigo}[title=app/src/main/AndroidManifest.xml]{XML}
<manifest xmlns:android=
    "http://schemas.android.com/apk/res/android">

  <uses-permission
      android:name="android.permission.INTERNET" />

  <application
      android:icon="@mipmap/ic_launcher"
      android:label="@string/app_name"
      android:theme="@style/Theme.LabMonitor">
    <activity
        android:name=".MainActivity"
        android:exported="true">
      <intent-filter>
        <action android:name="android.intent.action.MAIN" />
        <category android:name="android.intent.category.LAUNCHER" />
      </intent-filter>
    </activity>
  </application>
</manifest>
\end{codigo}
    \end{column}
    \begin{column}{0.40\textwidth}
      \small
      O manifesto declara ao sistema:
      \begin{itemize}
        \item as \textbf{permissões} necessárias;
        \item o \textbf{ícone}, o \textbf{nome} e o \textbf{tema};
        \item os \textbf{componentes} (Activities, Services\ldots);
        \item qual Activity abre pelo \textbf{ícone do launcher} (\cod{MAIN} + \cod{LAUNCHER}).
      \end{itemize}
      \begin{atencao}
        \cod{android:exported="true"} é obrigatório em componentes com \cod{intent-filter}.
      \end{atencao}
    \end{column}
  \end{columns}
\end{frame}

\begin{frame}[fragile]{MainActivity.kt: o ponto de entrada}
  \begin{columns}[T]
    \begin{column}{0.58\textwidth}
\begin{codigo}[title=MainActivity.kt]{Kotlin}
class MainActivity : ComponentActivity() {
  override fun onCreate(savedInstanceState: Bundle?) {
    super.onCreate(savedInstanceState)
    enableEdgeToEdge()  // desenha atrás das barras
    setContent {        // a partir daqui é Compose
      LabMonitorTheme {
        Scaffold(Modifier.fillMaxSize()) { padding ->
          Saudacao(
            nome = "Laboratório",
            modifier = Modifier.padding(padding)
          )
        }
      }
    }
  }
}

@Composable
fun Saudacao(nome: String, modifier: Modifier = Modifier) {
  Text(text = "Olá, $nome!", modifier = modifier)
}
\end{codigo}
    \end{column}
    \begin{column}{0.40\textwidth}
      \small
      \begin{itemize}
        \item \cod{ComponentActivity}: Activity base para Compose.
        \item \cod{onCreate}: chamado quando a tela é criada (Módulo 07).
        \item \cod{setContent \{\}}: define a interface com funções \cod{@Composable}.
        \item \cod{LabMonitorTheme}: tema Material 3 gerado em \arq{ui/theme/}.
        \item \cod{Scaffold}: estrutura básica de tela (barras, FAB).
      \end{itemize}
      \begin{dica}
        Os detalhes de Compose ficam para o Módulo 05. Por enquanto, rode e altere o texto.
      \end{dica}
    \end{column}
  \end{columns}
\end{frame}

% =============================================================================
\section{Gradle}

\begin{frame}{O que é o Gradle}
  \begin{columns}[c]
    \begin{column}{0.48\textwidth}
      \small
      \termo{Gradle} é o sistema de build: ele baixa dependências, compila, empacota, assina e roda testes.
      \begin{itemize}
        \item Os scripts são escritos em \textbf{Kotlin DSL} (\arq{.gradle.kts}).
        \item O \termo{Android Gradle Plugin (AGP)} ensina o Gradle a construir apps Android.
        \item \textbf{Gradle Sync}: a IDE lê os scripts e configura o projeto. Rode-o após qualquer mudança nos arquivos de build.
      \end{itemize}
    \end{column}
    \begin{column}{0.50\textwidth}
      \centering
      \begin{tikzpicture}[node distance=2.5mm]
        \node[caixaescura,minimum width=5.2cm] (set) {\arq{settings.gradle.kts}\\ \tiny repositórios e módulos incluídos};
        \node[caixa,minimum width=5.2cm,below=of set] (proj) {\arq{build.gradle.kts} (projeto)\\ \tiny plugins disponíveis para os módulos};
        \node[caixaambar,minimum width=5.2cm,below=of proj] (app) {\arq{app/build.gradle.kts}\\ \tiny configuração do app e dependências};
        \node[caixa,minimum width=5.2cm,below=of app,fill=cinzaclaro] (toml) {\arq{gradle/libs.versions.toml}\\ \tiny versões centralizadas (catálogo)};
        \draw[seta] (set) -- (proj);
        \draw[seta] (proj) -- (app);
        \draw[setatrac] (toml.east) -- ++(4mm,0) |- (app.east);
        \draw[setatrac] (toml.east) -- ++(4mm,0) |- (proj.east);
      \end{tikzpicture}
    \end{column}
  \end{columns}
\end{frame}

\begin{frame}[fragile]{Version catalog e dependências}
  \begin{columns}[T]
    \begin{column}{0.48\textwidth}
\begin{codigo}[title=gradle/libs.versions.toml,listing options={style=dm,language=bash,basicstyle=\ttfamily\tiny}]{bash}
[versions]
agp = "..."             # versões atuais
kotlin = "..."          # o assistente
composeBom = "..."      # preenche

[libraries]
androidx-compose-bom = { group = "androidx.compose",
  name = "compose-bom", version.ref = "composeBom" }
androidx-material3 = { group = "androidx.compose.material3",
  name = "material3" }

[plugins]
android-application = { id = "com.android.application",
  version.ref = "agp" }
\end{codigo}
    \end{column}
    \begin{column}{0.50\textwidth}
\begin{codigo}[title=app/build.gradle.kts,listing options={style=dm,language=Kotlin,basicstyle=\ttfamily\tiny}]{Kotlin}
plugins {
  alias(libs.plugins.android.application)
  alias(libs.plugins.kotlin.compose)
}
android {
  namespace = "br.edu.ifma.labmonitor"
  compileSdk = 36
  defaultConfig {
    applicationId = "br.edu.ifma.labmonitor"
    minSdk = 26
    targetSdk = 36
    versionCode = 1
    versionName = "1.0"
  }
}
dependencies {
  implementation(platform(libs.androidx.compose.bom))
  implementation(libs.androidx.material3)
  testImplementation(libs.junit)
}
\end{codigo}
    \end{column}
  \end{columns}
  \begin{caixadestaque}
    \textbf{BOM} (\emph{Bill of Materials}): fixa versões compatíveis de todas as bibliotecas Compose.
  \end{caixadestaque}
\end{frame}

\begin{frame}[fragile]{Configurações de dependência, tarefas e variantes}
  \begin{columns}[T]
    \begin{column}{0.5\textwidth}
      \footnotesize
      \begin{tabularx}{\linewidth}{@{}l L@{}}
        \toprule
        \cabecalho{Configuração} & \cabecalho{Uso} \\
        \midrule
        \cod{implementation} & código do app \\
        \cod{testImplementation} & testes unitários \\
        \cod{androidTestImplementation} & testes no aparelho \\
        \cod{debugImplementation} & só no build de debug \\
        \cod{ksp} & processadores de anotação \\
        \bottomrule
      \end{tabularx}
      \vspace{4pt}
      \begin{conceito}[Build types]
        \textbf{debug}: depurável, assinado com chave de teste. \textbf{release}: otimizado (R8), assinado com a sua chave, vai para a loja.
      \end{conceito}
    \end{column}
    \begin{column}{0.48\textwidth}
\begin{terminal}[title=Tarefas pelo terminal]
./gradlew tasks            # lista tarefas
./gradlew assembleDebug    # gera APK debug
./gradlew installDebug     # instala no aparelho
./gradlew test             # testes unitários
./gradlew lint             # análise estática
./gradlew bundleRelease    # gera AAB (loja)
./gradlew clean            # limpa o build
\end{terminal}
      \begin{dica}
        Os APKs gerados ficam em \arq{app/build/outputs/apk/}.
      \end{dica}
    \end{column}
  \end{columns}
\end{frame}

% =============================================================================
\section{Primeiro app e depuração}

\begin{frame}[fragile]{Compose Preview e execução}
  \begin{columns}[T]
    \begin{column}{0.52\textwidth}
\begin{codigo}[title=Preview]{Kotlin}
@Preview(showBackground = true, name = "Claro")
@Preview(
  uiMode = Configuration.UI_MODE_NIGHT_YES,
  name = "Escuro"
)
@Composable
fun SaudacaoPreview() {
  LabMonitorTheme {
    Saudacao("Bancada 3")
  }
}
\end{codigo}
      \begin{dica}
        O \emph{Preview} renderiza a função na IDE sem instalar o app. Use vários \cod{@Preview} para comparar tema claro e escuro, fontes grandes e telas diferentes.
      \end{dica}
    \end{column}
    \begin{column}{0.46\textwidth}
      \begin{conceito}[Executando]
        \begin{enumerate}
          \item Escolha o aparelho na barra superior.
          \item Clique em \textbf{Run} (\textcolor{verdeok}{$\blacktriangleright$}) ou \textbf{Shift+F10}.
          \item O Gradle compila, instala e abre o app.
        \end{enumerate}
      \end{conceito}
      \begin{conceito}[Live Edit]
        Com o app rodando, mudanças em funções Compose aparecem no aparelho sem reinstalar (\emph{Live Edit}).
      \end{conceito}
    \end{column}
  \end{columns}
\end{frame}

\begin{frame}[fragile]{Logcat e depurador}
  \begin{columns}[T]
    \begin{column}{0.5\textwidth}
\begin{codigo}[title=Registrando logs]{Kotlin}
private const val TAG = "LabMonitor"

fun lerSensor(valor: Double) {
  Log.d(TAG, "Leitura recebida: $valor")
  if (valor > 80.0) {
    Log.w(TAG, "Temperatura alta: $valor")
  }
}
\end{codigo}
      \footnotesize
      \begin{tabularx}{\linewidth}{@{}l l L@{}}
        \toprule
        \cod{Log.v} & verbose & detalhes extremos \\
        \cod{Log.d} & debug & desenvolvimento \\
        \cod{Log.i} & info & eventos normais \\
        \cod{Log.w} & warn & situação suspeita \\
        \cod{Log.e} & error & falhas \\
        \bottomrule
      \end{tabularx}
    \end{column}
    \begin{column}{0.48\textwidth}
      \begin{conceito}[Depurador]
        \begin{itemize}
          \item Clique na margem para criar um \emph{breakpoint}.
          \item Rode com \textbf{Debug} (ícone de inseto).
          \item Inspecione variáveis, avalie expressões e avance linha a linha (F8/F7).
        \end{itemize}
      \end{conceito}
      \begin{dica}[Filtro do Logcat]
        \cod{package:mine tag:LabMonitor level:debug} mostra só os logs do seu app.
      \end{dica}
      \begin{atencao}
        Nunca registre senhas, tokens ou dados pessoais em logs.
      \end{atencao}
    \end{column}
  \end{columns}
\end{frame}

\begin{frame}{Outras ferramentas da IDE}
  \small
  \begin{tabularx}{\textwidth}{@{}l L@{}}
    \toprule
    \cabecalho{Ferramenta} & \cabecalho{Para que serve} \\
    \midrule
    \textbf{Layout Inspector} & Ver a árvore de componentes do app em execução e contar recomposições \\
    \textbf{App Inspection} & Inspecionar o banco de dados (Room), o tráfego de rede e os \emph{workers} em tempo real \\
    \textbf{Profiler} & Medir CPU, memória e energia; achar travamentos (Módulo 14) \\
    \textbf{Device Explorer} & Navegar pelos arquivos do aparelho \\
    \textbf{Running Devices} & Espelhar e controlar o celular físico dentro da IDE \\
    \textbf{Gemini no Android Studio} & Gerar código, explicar erros e sugerir testes (sempre revise o resultado) \\
    \bottomrule
  \end{tabularx}
  \vspace{4pt}
  \begin{dica}[Atalhos que valem ouro]
    \textbf{Shift Shift}: busca global \quad \textbf{Alt+Enter}: correções rápidas \quad \textbf{Ctrl+Alt+L}: formatar \quad \textbf{Ctrl+B}: ir para a definição
  \end{dica}
\end{frame}

% =============================================================================
\section{Versionamento com Git}

\begin{frame}[fragile]{Git e GitHub no projeto}
  \begin{columns}[T]
    \begin{column}{0.52\textwidth}
\begin{terminal}[title=Fluxo básico]
git init
git add .
git commit -m "Projeto LabMonitor criado"
git branch -M main
git remote add origin https://github.com/<usuario>/LabMonitor.git
git push -u origin main
# a cada etapa:
git checkout -b feature/lista-equipamentos
git commit -am "Tela de lista de equipamentos"
\end{terminal}
    \end{column}
    \begin{column}{0.46\textwidth}
      \begin{conceito}[O que NÃO versionar]
        O \arq{.gitignore} gerado pelo Android Studio já exclui:
        \begin{itemize}
          \item \arq{build/}, \arq{.gradle/}, \arq{.idea/} (parcial);
          \item \arq{local.properties} (caminho do SDK na sua máquina).
        \end{itemize}
      \end{conceito}
      \begin{atencao}
        Nunca faça commit de \textbf{chaves de API}, \textbf{keystores} (\arq{.jks}) ou senhas. Se isso acontecer, considere a chave comprometida e gere outra.
      \end{atencao}
    \end{column}
  \end{columns}
\end{frame}

% =============================================================================
\section{Fechamento}

\begin{frame}{Resumo}
  \begin{resumo}
    \begin{itemize}
      \item Android Studio = editor + SDK + emulador + Gradle + ferramentas de inspeção.
      \item Rode apps no emulador (AVD) ou no celular real com depuração USB/Wi-Fi; o \cod{adb} faz a ponte.
      \item \arq{AndroidManifest.xml} declara componentes e permissões; \arq{MainActivity.kt} chama \cod{setContent}.
      \item O Gradle constrói o app; o \emph{version catalog} centraliza versões e o BOM alinha as bibliotecas Compose.
      \item \cod{minSdk}, \cod{targetSdk} e \cod{applicationId} ficam no \arq{app/build.gradle.kts}.
      \item Logcat, depurador e Layout Inspector são as ferramentas de diagnóstico do dia a dia.
    \end{itemize}
  \end{resumo}
\end{frame}

\begin{frame}{Glossário}
  \begin{glossario}
    SDK & conjunto de bibliotecas e ferramentas para desenvolver para uma plataforma \\
    IDE & ambiente integrado: editor, \emph{build}, execução e depuração \\
    Gradle & sistema de \emph{build} do Android, baseado em tarefas e \emph{plugins} \\
    Catálogo de versões & arquivo \cod{libs.versions.toml} que centraliza as dependências \\
    AVD & aparelho virtual configurado no emulador do Android \\
    ADB & ferramenta de linha de comando que conversa com aparelhos e emuladores \\
    Logcat & registro de mensagens do sistema e do app \\
    \emph{Breakpoint} & ponto em que o depurador pausa a execução \\
    Manifesto & \arq{AndroidManifest.xml}: componentes, permissões e metadados do app \\
  \end{glossario}
\end{frame}

\begin{frame}{Questões de revisão}
  \begin{questoes}
    \item Qual a diferença entre compilar com JIT e com AOT? Onde cada um aparece no Android?
    \item O que o Android SDK contém além das bibliotecas?
    \item Por que o Gradle consegue recompilar só parte do projeto?
    \item O que é uma dependência transitiva e por que ela pode causar conflitos?
    \item Por que o emulador precisa de aceleração por hardware?
    \item Cite duas situações em que o emulador não substitui um aparelho real.
    \item Para que servem o ADB e o Logcat?
    \item Descreva o método para depurar um app que fecha ao tocar num botão.
  \end{questoes}
\end{frame}

\begin{frame}{Exercícios}
  \begin{exercicio}[Exercício 1 --- Ambiente]
    Crie dois AVDs: um celular com a API mais recente e um com API 26. Rode o app nos dois e anote as diferenças visuais.
  \end{exercicio}
  \begin{exercicio}[Exercício 2 --- ADB]
    Pelo terminal, instale o APK de debug no seu celular, capture uma tela com \cod{adb shell screencap} e copie a imagem para o computador.
  \end{exercicio}
  \begin{exercicio}[Exercício 3 --- Gradle]
    Altere \cod{versionName} para \cod{"1.0-aula"}, gere o APK com \cod{./gradlew assembleDebug} e localize o arquivo gerado. O que acontece se \cod{minSdk} for maior que a API do seu emulador?
  \end{exercicio}
  \begin{projeto}[LabMonitor --- Etapa 1]
    Crie o projeto LabMonitor (\cod{br.edu.ifma.labmonitor}), troque o texto inicial por ``LabMonitor --- DEE/IFMA'', adicione um \cod{Log.i} no \cod{onCreate} e publique o repositório no GitHub.
  \end{projeto}
\end{frame}

\begin{frame}{Para aprofundar}
  \small
  \begin{itemize}
    \item Conheça o Android Studio: \url{https://developer.android.com/studio/intro}
    \item Rodar apps em aparelho físico: \url{https://developer.android.com/studio/run/device}
    \item Emulador: \url{https://developer.android.com/studio/run/emulator}
    \item ADB: \url{https://developer.android.com/tools/adb}
    \item Configurar o build: \url{https://developer.android.com/build}
    \item Version catalogs: \url{https://developer.android.com/build/migrate-to-catalogs}
  \end{itemize}
\end{frame}

\begin{frame}{Respostas das questões de revisão}
  \begin{respostas}
    \item JIT compila durante a execução; AOT, antes. O ART combina os dois: compila AOT partes do app (com perfis) e usa JIT no restante.
    \item Ferramentas de \emph{build} e de linha de comando (ADB), imagens de sistema para o emulador, documentação e as plataformas de cada versão da API.
    \item Porque organiza o trabalho em tarefas com entradas e saídas conhecidas: se as entradas não mudaram, reaproveita o resultado (\emph{build} incremental e cache).
    \item É a dependência de uma dependência. Duas bibliotecas podem exigir versões diferentes da mesma biblioteca; o gerenciador precisa escolher uma compatível.
    \item Porque executa um sistema completo numa máquina virtual; com VT-x, AMD-V ou KVM as instruções rodam direto no processador.
    \item Uso de Bluetooth, NFC ou USB; medição de desempenho e bateria; câmera e sensores reais.
    \item O ADB instala, inicia e depura apps e abre um terminal no aparelho; o Logcat mostra as mensagens e a pilha de erros.
    \item Reproduzir o fechamento, ler a pilha no Logcat, colocar um \emph{breakpoint} antes da linha suspeita, testar a hipótese, corrigir e criar um teste.
  \end{respostas}
\end{frame}

\encerramento{Módulo 02 --- Kotlin essencial: tipos, null safety, funções, classes e coleções.}

\end{document}
